병렬 이분 탐색(parallel binary search)은 매개변수 탐색 과정을 여러 번 처리해야 할 때 빠르게 처리하기 위한 방법이다. 기본적으로 매개변수 탐색을 여러 번 처리할 때 같은 값 $x$에 대해 $x$가 답이 될 수 있는지 판단하는 것이 중복적으로 나타나고, 이 시간을 줄이겠다는 아이디어에서 만들어진 알고리즘이다. 이번 장에서는 6장에서 배운 매개변수 탐색을 바탕으로 병렬 이분 탐색을 하는 방법을 배운다.

\textbf{* 주의: 병렬 이분 탐색은 뒤쪽 챕터의 알고리즘과 함께 쓰이는 경우가 전부이므로, 뒤의 내용을 공부하다가 필요한 시점에 다시 보는 것을 추천한다.}

\section{쿼리와 이분 탐색}
\Cref{ex: parametric-search}에 쿼리를 추가한 문제를 생각하자.
\begin{example} \label{ex: parallel-binary-search}
    자연수로만 이루어진 어떤 배열 $A$가 주어질 때, 다음 쿼리 $Q$개를 처리하는 프로그램을 작성하라.
    \begin{itemize}
        \item 쿼리: $A_i$ 중 $x$보다 큰 수의 개수가 $k$개 이하가 되는 $x$의 최솟값을 구하여라.
    \end{itemize}
\end{example}

만약 매개변수 탐색을 $Q$번 진행한다면 시간복잡도가 $O(NQ\log X)$가 될 것이다. 하지만, 병렬 이분 탐색을 사용할 경우 $O(N \log Q \log X + Q \log X)$로 줄일 수 있다.

각 쿼리의 $k$를 $q_i$라고 놓자. 그리고 각 쿼리마다 현재 가능한 값의 범위 $[s_i, e_i]$를 생각하자. 초깃값은 모든 $i$에서 $s_i=1, e_i=X$일 것이다. 이제 $x=m = (s+e)/2$에서 $x$보다 큰 수의 개수를 세자. 이 개수보다 $q_i$가 큰 쿼리는 $s_i$가 바뀌고, $q_i$가 그 개수 이하인 쿼리는 $e_i$가 바뀔 것이다.

이제 각 $[s_i, e_i]$가 절반으로 감소하였다. 이와 동시에 가능한 $x=m$의 값은 두 가지가 되었다. 만약 이 두 가지 $x=m$에 대해 각각 $O(N)$의 시간으로 $x$보다 큰 수의 개수를 계산한다면 시간복잡도의 개선이 전혀 일어나지 않을 것이다(\Cref{pb: wrong-pbs-time}). 병렬 이분 탐색을 통해 시간복잡도를 개선하려면 \textbf{한 번의 탐색으로 다양한 $x$ 값에 대해 원하는 값(여기서는 $x$보다 큰 수)}을 구할 수 있어야 한다.

이를 위해 이분탐색을 사용할 수 있다. 가능한 각 쿼리에서 $m$의 값을 $x_1, x_2, \cdots, x_q$라고 놓고 이를 정렬하자. 그리고 배열 $c_i$에 값이 $x_{i-1}$ 초과 $x_i$ 이하인 수의 개수를 저장하자(편의상 $x_0=0$, $x_{q+1}=X$로 놓는다). 이 과정은 $A_i$ 초과인 가장 앞에 있는 수 $x_j$를 배열 $x$에서 이분탐색을 통해 찾으면, $(x_{j-1}, x_j]$에 $A_i$가 존재하므로 $c_j$를 $1$ 증가시키면 된다. 따라서 $O(N \log Q)$의 시간에 $c_i$를 계산할 수 있다. 또한 $c_i$의 정의에 의해 $i+1$부터 $q+1$까지의 $c_i$의 부분 합이 각 $x_i$보다 큰 수의 개수이므로 각 $x_q$에 대한 답을 $O(Q)$의 시간에 찾을 수 있다. 구현은 다음 \Cref{code: parallel-binary-search}을 참고하라.

\begin{code} \label{code: parallel-binary-search}
\Cref{ex: parallel-binary-search}의 정답 코드
\begin{lstlisting}
#include <bits/stdc++.h>
using namespace std;

int n, arr[101010], q, k[101010];
int s[101010], e[101010], c[101010], sum[101010];
const int X = (1<<30);

int main() {
    cin >> n;
    for (int i = 1; i <= n; i++) cin >> arr[i];
    cin >> q;
    for (int i = 1; i <= q; i++) cin >> k[i];

    for (int i = 1; i <= q; i++) {
        s[i] = 1;
        e[i] = X;
    }
    while (s[1] < e[1]) {
        for (int i = 1; i <= q + 1; i++) c[i] = sum[i] = 0;
        vector< pair<int, int> > m;
        m.push_back({0, 0});
        for (int i = 1; i <= q; i++) m.push_back({(s[i] + e[i]) / 2, i});
        m.push_back({X, q + 1});
        sort(m.begin(), m.end());

        for (int i = 1; i <= n; i++) {
            int j = lower_bound(m.begin(), m.end(), make_pair(arr[i] + 1, 0)) - m.begin();
            c[j]++;
        }

        for (int i = 1; i <= q + 1; i++) sum[i] = sum[i - 1] + c[i];
        for (int i = 1; i <= q; i++) {
            if (sum[q + 1] - sum[i] <= k[m[i].second]) e[m[i].second] = m[i].first;
            else s[m[i].second] = m[i].first + 1;
        }
    }

    for (int i = 1; i <= q; i++) cout << s[i] << " ";

    return 0;
}
\end{lstlisting}
\end{code}

정리하면 병렬 이분 탐색은 최적해 문제 쿼리를 업데이트 쿼리와 결정 문제 쿼리로 환원하는 것이다:
\begin{itemize}
    \item 업데이트 쿼리: 결정 문제의 고정된 기준값 $m$이 증가할 때, 결정 문제를 빠르게 해결하기 위해 저장해야 하는 값들을 갱신
    \item 결정 문제 쿼리: 주어진 $m$에 대해 결정 문제를 해결
\end{itemize}

이처럼 업데이트 쿼리가 추가되기 때문에 업데이트가 있는 상황에서 주어진 문제를 빠르게 해결하기 위해 사용되는 자료구조들과 함께 쓰이는 경우가 많다.

\section{연습문제 A}
\begin{problem} \label{pb: wrong-pbs-time}
    \Cref{ex: parallel-binary-search}에 대한 병렬 이분 탐색 과정에서 $Q$개의 쿼리에서 나온 모든 $m$에 대해 $O(N)$의 시간으로 탐색한다면 시간복잡도가 어떻게 되는가?
\end{problem}
\begin{problem}
    \Cref{code: parallel-binary-search}에서 $X$를 \texttt{1<<30}으로 놓은 이유가 무엇인가?
\end{problem}
\begin{problem}
    \Cref{code: parallel-binary-search}에서 부분 합을 구간 $[i+1, q+1]$로 처리한 이유가 무엇인가? 왜 $[i+1, q]$가 아닌가?
\end{problem}
\begin{problem}
    사실 \Cref{ex: parallel-binary-search}은 병렬 이분 탐색 없이 해결할 수 있다. 병렬 이분 탐색 없이 해결하는 코드를 짜고, 그 코드의 시간복잡도를 구하라.
\end{problem}

\section{연습문제 B}

\begin{problem}
    $x_i$가 $x_{i-1} \cdot a + b$를 $10^9+7$로 나눈 나머지인 점화관계를 가진 길이 $N$의 수열 $x_i$에 대해 $x_i$에서 $q$번째로 작은 원소를 찾는 쿼리 $Q$개를 처리하는 프로그램을 작성하라. 정확히 말하면 $y$보다 작은 $x_i$의 원소의 개수가 $q$인 가장 작은 $y$를 찾아야 한다. \textbf{단, $x_i$를 모두 저장할만큼 메모리가 충분하지 않다.}
    \begin{itemize}
        \item (12926번) $N \leq 10^6$, $Q \leq 100$, $x_0, a, b$는 $10^9+7$ 미만의 음 아닌 정수
        \item (12921번) $N \leq 2 \cdot 10^6$, $Q \leq 1000$, $x_0, a, b$는 $10^9+7$ 미만의 음 아닌 정수
    \end{itemize}
\end{problem}

\section{연습문제 C}
\textbf{* 주의: 본 챕터의 연습문제 C는 뒤에 나오는 자료구조를 알아야 풀 수 있는 문제로 구성되어 있다.} 만약, 이 책을 처음부터 읽고 있으며 뒤의 자료구조에 관한 지식이 없다면 연습문제 C를 보며 아이디어만 생각하거나 나중에 다시 풀어보자.

\begin{problem}
    (8217번) 원이 $M$ 분할되어 각각 $1$번부터 $M$번까지 번호가 붙어 있다. 특히, $1$번과 $M$번은 붙어 있다. 그리고 각각은 $N$명의 사람 중 한 명이 소유하고 있으며 이 정보는 입력으로 주어진다. 또한 $Q$일 동안 연속된 부채꼴 $[l_i, r_i]$번에서 각각 $a_i$만큼의 석유가 나온다($1 \leq i \leq Q$). 참고로 $l_i > r_i$이면 $[l_i, r_i]$는 $[l_i, m] \cup [1, r_i]$를 뜻한다. 각 사람이 자신이 가진 땅의 석유를 매일 모을 때, 각 사람마다 원하는 목표치인 $p_j~(1 \leq j \leq N)$ 이상의 석유를 얻기 위해 기다려야 하는 날의 수를 출력하라. $N \leq 3 \times 10^5$, $M \leq 3 \times 10^5$, $Q \leq 3 \times 10^5$이고 $a_i \leq 10^9$, $p_i \leq 10^9$이다.
\end{problem}

\begin{problem}
    (9877번) $10^9$ 이하의 음 아닌 정수가 적힌 $N \times M$ 격자에서 정수 $T$가 주어질 때 각 셀에 대해 각 칸에 대해 다음 쿼리를 처리하라. $N \leq 500$, $M \leq 500$이다.
    \begin{itemize}
        \item 상하좌우로 인접한 칸 중 적힌 수의 차이가 $D$ 이하인 칸으로 이동하는 것만 허용할 때 시작 칸 $(a, b)$에서 도달할 수 있는 칸의 개수가 $T$개 이상이 되도록 $D$를 정하자. 가능한 $D$의 최솟값을 출력하라.
    \end{itemize}
\end{problem}

\begin{problem}
    (13952번) 방해물이 있는 $N \times N$ 격자에서 다음 쿼리 $Q$개를 처리하라. $N \leq 1000$, $Q \leq 3 \times 10^5$이다.
    \begin{itemize}
        \item $(x_a, y_a)$가 중심인 홀수 길이의 정사각형이 방해물에 걸리지 않고 상하좌우로 평행하게 이동하여 중심이 $(x_b, y_b)$인 위치까지 이동할 수 있는 정사각형 길이의 최댓값을 출력하라. 그 어떤 정사각형도 불가능하면 $0$을 출력하면 된다.
    \end{itemize}
\end{problem}

\begin{problem}
    (16976번) 컴퓨터를 한 대 갖고 있는 $N$명의 사람이 $M$종류의 게임 중 하나의 게임을 구매했다. 그리고 순차적으로 $Q$번 두 사람이 가진 컴퓨터 사이에 랜선을 연결할 것이다. $i$번 게임을 구매한 모든 사람의 컴퓨터가 연결되었을 때 $i$번 게임을 시작할 때, 각 게임이 몇 개의 랜선을 연결한 이후에 시작될 수 있는지 출력하라. $N \leq 10^5$, $M \leq 10^5$, $Q \leq 10^5$이다.
\end{problem}

\begin{problem}
    (20654번) $N$개의 액체가 있고, 각 액체는 맛 $d_i$와 단위용량당 가격 $p_i$가 정해져 있다. 이 액체 중 몇 개를 원하는 양만큼 넣어 음료를 만들고 싶다. 단, $i$번 액체는 $l_i$ 이하만큼 넣어야 한다. 이때 다음 쿼리 $M$개를 처리하는 프로그램을 작성하라. $N \leq 10^5$, $M \leq 10^5$이고 $d_i, p_i, l_i$는 $10^5$ 이하의 자연수, $g_i, L_i$는 $10^{18}$ 이하의 자연수이다.
    \begin{itemize}
        \item 음료를 만드는데 필요한 액체의 가격이 $g_j$ 이하이며 음료의 용량이 $L_j$ 이상인 음료 중에서 음료에 포함되는 액체의 맛 $d_i$의 최솟값의 가능한 최댓값을 출력하라.
    \end{itemize}
\end{problem}

\begin{problem}
    (21095번*) 길이 $N$의 수열 $A_i$가 주어질 때, 다음 쿼리 $Q$개를 처리하는 프로그램을 작성하라.
    \begin{itemize}
        \item $l, r, k$가 주어질 때 모든 $1 \leq i \leq k$에 대해 $C_i + C_{(i~\textrm{mod}~k)+1} \leq W$이면서 $A[l..r]$에서 몇 개의 수를 지워 $k$개의 수를 남긴 부분 수열인 $C_i$가 존재하는 $W$의 최솟값을 출력하라.
    \end{itemize}
\end{problem}

\begin{problem}
    (25458번*) \url{https://www.acmicpc.net/problem/25458}
\end{problem}

\begin{problem}
    (27199번*) 주어진 홀수 길이의 배열에서 \textbf{연속된 세 숫자 중 중앙값만 남기는 연산}을 배열의 길이가 $1$이 될 때까지 해 남길 수 있는 수의 최댓값을 배열의 \textit{중앙값 평균}이라고 하자. 길이 $N$의 배열 $A_i$와 $Q$개의 $(l, r)$ 쌍이 주어질 때 $A[l..r]$의 \textit{중앙값 평균}을 출력하는 프로그램을 작성하라. $N \leq 50000$, $Q \leq 10^5$, $0 \leq A_i \leq 10^9$이고 모든 쿼리에서 $r-l+1$은 홀수이다.
\end{problem}